\section{自己開示の場としてのカラオケ}\label{自己開示の場としてのカラオケ}
カラオケは，音楽を介した相互作用が生まれる社会的な場である．本節では，カラオケにおける自己開示と受容の定義を明確にし，その特徴について論じる．

\subsection{カラオケにおける自己開示と受容}
\subsection*{カラオケにおける自己開示の定義}

自己開示とは，\ref{自己開示の定義}節でも述べたように，言語的なコミュニケーションに限らず，非言語的な表現も一形態である．また，受容についても非言語的・言語的の両面から評価可能な形で定義づける必要がある．

単に曲名を相手に伝えるだけでは，自己開示にならない．そこで，曲を実際に流し，歌唱をする，またそれに伴いなぜそれが好きであるか，なぜそれを選曲したかといった動機を相手に伝える行為こそが，カラオケにおける自己開示になると考える．

そこで，本研究における自己開示とは，音楽選択による自己表現と，その選択に関する言語的説明の2つの要素とする．音楽選択による自己表現とは，個人の好みや価値観が反映された楽曲を選択し他者と共有する行為を指す．選択に関する言語的説明とは，選曲理由や楽曲にまつわる個人的経験，感情などの言語的な共有を指す．


\subsection*{カラオケにおける受容の定義}
本研究における受容とは，①相手の選択した音楽を傾聴する非言語的受容と，②音楽や説明に対する共感や理解を示す言語的受容の2つの要素から構成される．これらの要素は，相手の自己開示に対する肯定的な反応として機能する．


\subsection{カラオケにおける自己開示の特徴}

カラオケにおける自己開示の特徴として，実体験を伴う開示が挙げられる．通常の会話では音楽の好みを口頭で開示することに留まり，曲名や好みの表明のみでは相手に音楽の本質的な魅力を伝えることは難しい．一方でカラオケでは，実際の楽曲が流れ，その曲調やメロディー，リズムといった音楽的特徴を直接体験として共有できる．さらに，歌唱を通じて自身がその楽曲のどの要素に魅力を感じているのかを，より具体的かつ直接的に表現することが可能となる．
また，カラオケにおける自己開示は，好みの曲の共有から感情の開示まで，段階的な開示が存在すると考えられる．好みのアーティストの楽曲を選択することから始まり，その選曲理由や曲への思い入れが語られ，さらに歌唱を重ねる中で思い出や個人的な感情が共有される．このように，カラオケは音楽の選曲と歌唱を通じて，徐々に内面を開示していくためのきっかけとして機能すると考えられる．

さらに重要な特徴として，明確な役割分担による相互性の確保がある．数家\cite{数家鉄治_2008}が指摘するように，カラオケは共同行為に必要な自己と他者の相互依存性をより緊密にできる場である．「歌を歌う」「相手の歌を聴く」という明確な役割が交互に繰り返される構造は，自己開示の相互性を自然な形で実現する．これにより通常の会話場面では，実現が困難である明示的な相互の自己開示や，自己開示のバランスの制限をすることが可能となる．

\vskip\baselineskip

これらの特徴から，カラオケが音楽嗜好の共有を通じた自己開示の環境・状況として有効であると考える．